\section{Úvod}

Princípy čístého kódu, ako napríklad malé funkcie, zmysluplné pomenovania, minimalizácia komentárov a princíp DRY (Don't Repeat Yourself), sa stali súčasťou štandardnej výučby programovania a sú často považované za správny prístup k vývoju softvéru vo veľkom množstve spoločností a boli široko prijaté agilnými tímami po celom svete \cite{CLEANCODE_MARTIN2008}.

Pôvodným cieľom Clean Code bolo zvýšiť čitateľnosť, udržiavateľnosť a kvalitu softvéru. V priebehu rokov sa však objavili názory, že niektoré odporúčania a postupy boli prijaté príliš rýchlo a nie sú nutne overené komunitou programátorov. Existujú názory, že snaha o dokonalý kód môže viesť k predčasnej optimalizácii, zbytočnému pridávaniu funkcionalít a nadmernému navrhovaniu systémov (over-engineeringu).

Tento problém je obzvlášť aktuálny v prostredí agilného vývoja, kde je prioritou rýchle dodávanie hodnoty používateľovi. Silná orientácia na architektonickú čistotu môže spomaľovať vývoj, zvyšovať komplexnosť systému a odvádzať pozornosť od skutočných potrieb používateľov. Výskumná otázka tejto práce teda znie:

\textbf{Vedie dogmatické uplatňovanie princípov Clean Code k predčasnej optimalizácii, preťaženiu funkcionalitami a over-engineeringu?}

Na zodpovedanie tejto otázky bude analyzovaných zopár názorov odborníkov z praxe, akademických zdrojov a súčasných diskusií o softvérovom inžinierstve. Cieľom je identifikovať situácie, v ktorých Clean Code prináša hodnotu, a situácie, v ktorých sa môže stať kontraproduktívnym.